\chapter{绪论}

\section{研究背景及意义}
人类大脑作为自然界最复杂的信息处理系统，包含约860亿个神经元与海量的突触连接，承载着思维、感知、记忆、情绪、决策等高级认知功能\cite{herculano2009human}。大脑神经网络通过复杂的电化学活动产生丰富的神经电信号，这些信号以不同的时空尺度反映着神经元群体的协同活动模式，构成了认知过程的直接电生理表现。脑电图（Electroencephalography, EEG）是通过头皮电极记录由同步突触后电位汇聚而成的微弱电信号，反映了大脑皮层大规模神经元群体活动的时变特性，具有毫秒级时间分辨率、非侵入、安全与成本低等优势\cite{nunez2006electric}。自Caton首次观察到动物大脑的电活动\cite{caton1970electric}，到Berger记录到人类EEG并发现$\alpha$节律\cite{berger1930ueber}，EEG已成为认知神经科学与脑机接口（Brain--Computer Interface, BCI）研究的核心数据来源。

\begin{figure}[htbp] 
  \centering
  \includegraphics[width=\textwidth]{C1-1BCI组成.pdf}
  \caption{EEG-BCI系统的主要组成\cite{gu2021eeg}}
  \label{fig:C1-1}
\end{figure}

如图~\ref{fig:C1-1}所示，基于EEG的BCI（Electroencephalography-based Brain-Computer Interface, EEG-BCI）系统旨在建立大脑活动与外部设备间的直接通信，使用户绕过神经--肌肉通路，通过意图或注意的改变控制外部系统。如图~\ref{fig:C1-2}所示，EEG-BCI系统已在辅助交流与康复训练、轮椅与机械臂控制、游戏交互与注意力训练，以及情绪与疲劳监测等场景中展现出应用潜力。
现有EEG-BCI广泛采用三类经典范式：运动想象（Motor Imagery, MI）利用感觉运动皮层的事件相关去同步化（Event-Related Desynchronization，ERD）和事件相关同步化（Event-Related Synchronization，ERS）特性实现左/右手、双脚等意向识别\cite{pfurtscheller1999event}；P300拼写范式基于注意驱动的事件相关电位（Event-Related Potential, ERP），在oddball刺激中目标出现约300毫秒后在顶叶区出现正向电位\cite{farwell1988talking}；稳态视觉诱发电位（Steady-State Visual Evoked Potential, SSVEP）范式通过识别枕叶皮层对不同频率闪烁刺激的同步响应实现快速选择\cite{bin2011high}。这些范式在受控环境下取得了较高准确率和信息传输率。



\begin{figure}[htbp] 
  \centering
  \includegraphics[width=0.55\textwidth]{C1-2BCI应用.pdf}
  \caption{EEG-BCI系统典型应用场景}
  \label{fig:C1-2}
\end{figure}

现有EEG-BCI系统面临两类核心瓶颈：
第一，信噪比低且非平稳。微弱的脑皮层电信号极易淹没在强大的生理噪声（如眼动、肌电、心电伪迹）和环境电磁干扰背景中\cite{goncharova2003emg}。在实际应用场景下，用户不可避免的头部移动、眨眼、面部肌肉活动，以及环境中的电子设备干扰，都会显著劣化信号质量。这种信噪比的大幅波动和下降，直接导致特征提取困难、意图解码错误率提升，使得系统性能不稳定甚至完全失效。信号质量的恶化不仅影响系统的实时性能，还会显著增加用户的认知负担，用户需要在噪声环境中更加专注地产生清晰的神经信号，这种额外的努力往往会导致疲劳和挫败感，进而影响用户的使用体验和系统的长期可用性；
第二，信号标注依赖极强。监督学习依赖高质量信号片段和标签配对的数据进行训练，而获取这类高质量标注数据需要在严格控制的实验环境中进行冗长的数据采集过程\cite{lotte2018review}。被试需要重复执行特定的认知任务，同时研究人员记录每个时刻的任务状态和神经响应，这一过程不仅成本高昂、效率低下，而且获得的数据往往局限于特定的实验条件和任务背景，泛化能力有限。同时，大量的日常神经电活动数据由于缺乏精确的任务标签而无法被有效利用，造成了数据资源的浪费。这种标注数据的稀缺性和获取难度导致了BCI系统在面对新用户、新任务或新环境时往往需要重新进行大量的标定工作，严重制约了系统的部署效率和实用性。此外，标注过程本身的主观性和不一致性也会引入额外的噪声，影响模型的训练效果和泛化能力。


与此同时，人工智能领域正经历从判别式模型向生成式模型的深刻范式转变\cite{ooi2025potential}。以变分自编码器（Variational Autoencoder, VAE）\cite{kingma2014vae}、生成对抗网络（Generative Adversarial Network, GAN）\cite{goodfellow2014generative}、正规化流\cite{rezende2015variational}与扩散模型\cite{ho2020denoising,song2021score}为代表的生成建模方法，能够学习数据分布并合成高保真样本；如图~\ref{fig:C1-3}所示，以大规模预训练与自监督为特征的基础模型与多模态模型（如CLIP\cite{radford2021learning}、GPT\cite{brown2020language}等）利用无标注数据学习通用表征并在开放任务中泛化\cite{rombach2022high}，展现出了创造复杂、多样、高质量内容的惊人能力 \cite{chavez2023chat, ramesh2021zero, rombach2022high}。
生成式模型通过概率性生成机制和多样性采样策略，能够在输入存在噪声或不完整的情况下仍然产生合理的输出。其普遍采用自监督学习策略，能够从大量无标注数据中学习丰富的数据表征和内在规律。例如，GPT模型通过预测下一个词的自监督任务从互联网文本中学习语言知识，扩散模型通过去噪任务从图像数据中学习视觉表征，这种学习方式显著降低了对人工标注数据的依赖，使得模型能够利用互联网超大规模的数据进行训练。因此，生成式模型能够在连续的高维空间中生成无穷多样的输出内容，从根本上突破了表达能力的界限。
生成式建模的鲁棒生成机制、自监督表征学习与开放表达空间的优势，与EEG-BCI的两大瓶颈形成高度互补，为重塑EEG的建模范式提供了坚实路径。

\begin{figure}[htbp] % use float package if you want it here
  \centering
  \includegraphics[width=\textwidth]{C1-3AIGC应用.pdf}
  \caption{生成式人工智能的主要应用\cite{zhang2025generative}}
  \label{fig:C1-3}
\end{figure}


将生成式建模应用于EEG-BCI的思路突破了将EEG视为待分类的静态特征向量的传统认知，转而将其概念化为具有内在生成规律的动态系统观测输出，通过学习和建模神经活动背后的生成机制和潜在结构，从概率分布的角度理解神经系统的变异性和不确定性。在学习策略层面，生成式建模实现了从依赖外部监督标签的被动学习向基于数据内在结构的主动发现转变，通过自监督学习策略从EEG的时空结构和频谱特性等内在模式中自主发掘通用的特征表示，为跨试次、跨被试和跨任务的知识迁移奠定基础。
在应用模式层面，生成式建模将EEG处理从问题特定的分类工具提升为通用的认知建模接口，不仅能够识别已知的神经模式，更具备生成新的神经活动模式的创造能力，在信号重建、特征表示、数据增强等多个层面开启了广泛的应用可能性。例如，其能够处理不完整、缺失或损坏的EEG数据，通过学习到的分布知识来填补缺失信息，重建高质量的信号。这种能力对于提高便携式BCI设备的实用性具有重要意义，因为这类设备往往受到电极数量限制、功耗约束和环境干扰等因素的影响，导致采集到的信号质量相对较低。
然而，多源EEG的异构性、非平稳性和跨尺度动态特性为其生成式建模方法带来了较大的技术挑战。数据来源包含不同采集设备、实验条件、被试个体和时间状态的多重变异源，神经活动的时空多尺度特性要求模型同时捕获从微秒级神经元同步到分钟级适应性变化的跨尺度动态模式，分布式编码特性和高维相关性结构使得标准的独立同分布假设不再适用。因此，需要设计多分辨率特征提取、动态建模等新技术来构建适应EEG特性的生成式建模方案。

本课题以国家科技创新2030--“脑科学与类脑研究”重大项目《侵入式/非侵入式双向闭环脑机接口关键技术及系统》为背景，面向多源EEG数据，探索从判别式分类到生成式理解的范式转变。围绕信号级重建与增强、通用表征提取等关键问题，本研究提出双层级生成式建模框架，以概率建模刻画EEG的生成机制和潜在结构，以自监督方式学习可迁移的通用表征，并在典型的EEG-BCI任务中进行验证。最后，本研究构建了轻量化可穿戴硬件平台，并开发了生成式算法驱动的认知状态实时监测与调节的系统原型。

\section{脑电信号分析研究现状}
\subsection{脑电信号的基础特性}
EEG是通过头皮电极记录的电位时间序列，主要反映皮层锥体细胞同步突触后电位的体表叠加效应，具有毫秒级时间分辨率与良好的安全性。如图~\ref{fig:C1-4}所示，国际标准化电极放置方案包括10--20系统、10--10系统等多种配置，常见通道数为8--256，采样率多为250--1000 Hz\cite{klem1999ten}。本节从信号生成与锁时特性、频谱，以及拓扑结构等方面阐述EEG特性。

\begin{figure}[htbp] % use float package if you want it here
  \centering
  \includegraphics[width=0.97\textwidth]{C1-4国际1010.pdf}
  \caption{国际10--10系统电极放置方案\cite{klem1999ten}}
  \label{fig:C1-4}
\end{figure}

从生成机制与锁时特性来看，EEG 可沿两条正交轴加以刻画：一是内源—外源，指活动是否依赖离散的外部/行为事件；二是相位锁定—诱导，指相位是否随事件稳定对齐。内源活动不依赖瞬时外部事件，常在静息或持续任务中自发出现；外源活动由离散刺激、动作启动等事件诱发。相位锁定成分对事件具有稳定相位，跨试次平均后仍能保真；诱导成分主要体现为功率或同步性的变化，其相位不随事件稳定对齐。在该框架下，典型成分可归纳如下：

（1）外源+相位锁定：ERP如P1/N1（感觉加工）、N2（冲突/抑制）与 P3（目标检测与情境更新），因相位稳定，适于通过跨试次叠加平均提高信噪比；SSVEP是相位锁定的极端情形，即在恒定频率的视觉调制下，皮层在刺激频率及其谐波处出现显著谱峰，单试次即可可靠检测；运动相关皮层电位（Movement-Related Cortical Potential, MRCP）在运动启动前约1s出现缓慢负位移，时间上可稳定锁定于动作启动事件；

（2）外源+诱导：ERS/ERD表现为特定频带相对基线的功率增强或抑制。感觉运动节律（Sensory Motor Rhythm, SMR）在运动准备、执行或动作想象时通常呈ERD，运动终止后常出现$\beta$反弹型ERS；

（3）内源+诱导：缓慢皮层电位（Slow Cortical Potential，SCP）跨数秒乃至更久，反映皮层兴奋性基线的重设，通常不锁定离散事件，可在神经反馈训练中被意志性调控；自发节律（如$\alpha/\mu/\beta$）不锁定具体事件，但其功率与相位统计随任务需求与警觉水平系统性变化，是认知状态表征的重要载体。


从频谱维度与功能关联看，EEG的功率分布近似呈$1/f$衰减，叠加多条功能相关节律带。表~\ref{tab:bands}汇集了常用频段、频率范围、典型功能关联及代表性BCI应用。这些频段边界更多是研究上的约定先验，实际个体的峰频、峰拓扑与调制强度存在差异。例如，枕区$\alpha$在闭眼静息明显增强，视觉与注意任务通常引发抑制；位于中央沟两侧的感觉—运动区$\mu/\beta$节律对运动准备与执行尤其敏感，是运动想象BCI的核心表征；高于30 Hz的$\gamma$及高伽马活动（>70 Hz）在头皮EEG中受背景噪声与肌电干扰影响更大，但在高质量采集与特定范式下仍可作为感知绑定与精细加工的指示信号。

\begin{table}[htbp]
\centering
\caption{常用的EEG频段与功能关联}
\label{tab:bands}
\begin{tabular}{lccc}
\hline
频段 & 频率范围（Hz） & 典型功能关联 & 代表性BCI应用 \\
\hline
$\delta$ & 1--4 & 深睡、病理状态 & 意识评估 \\
$\theta$ & 4--8 & 记忆编码、工作负荷 & 工作负荷监测 \\
$\alpha$ & 8--13 & 闭眼静息、注意抑制 & MI-ERD/ERS判别 \\
$\beta$ & 13--30 & 运动准备、感觉-运动整合 & MI-ERD/ERS判别 \\
$\gamma$ & 30--50（或更高） & 感知绑定、注意 & 认知状态估计 \\
\hline
\end{tabular}
\end{table}

在空间拓扑与网络组织上，EEG呈现出既局部又分布式的混合特性。一方面，体积传导与颅骨高阻抗导致头皮电位在空间上的扩散，使单源的场型在多电极上可见，从而削弱了空间定位的特异性；另一方面，不同节律在不同皮层区具有相对稳定的优势分布，如枕区$\alpha$、中央区$\mu/\beta$、顶区$\theta$，这使得恰当的参考选择与空间滤波可以显著提高局部信噪比并增强判别力。进一步地，跨区相互作用与耦合（如相干性、相位锁定值、加权相位滞后指数，以及相位振幅耦合）揭示了网络层面的同步与信息整合，是刻画任务驱动重构与病理网络改变的关键维度。介于局部与网络之间的微状态分析，则以数十至上百毫秒的稳定电位拓扑序列为基本单元，描述静息与持续任务中的动态脑态切换，有助于从全头皮视角理解大脑网络状态的切换规律和时空组织模式。

这些复杂的时空特性使得EEG信号的分析和建模面临独特的挑战。例如，时空多尺度要求EEG分析方法能够同时捕获从毫秒级的瞬时响应到秒级的状态变化，个体差异则要求EEG分析模型具备良好的泛化和适应能力。传统的信号处理方法往往难以同时兼顾这些特性，而生成式建模方法通过学习数据的内在分布和生成机制，为处理这些复杂特性提供了新的可能性。

\subsection{基于特征工程的脑电分析方法}
EEG分析方法经过数十年发展，形成了一套经典的分析流程和技术体系，主要包括预处理、特征提取和选择，以及分类识别三个关键环节。这些方法为理解神经活动机制和开发早期BCI系统奠定了基础。本节将系统阐述这些传统方法的技术原理、应用场景及其局限性。

（一）信号预处理方法

EEG预处理旨在提高信号质量，消除采集过程中的各类噪声和伪迹，为后续分析提供高信噪比的数据基础。根据处理对象和技术原理的不同，预处理技术主要包括空间滤波、时域滤波、伪迹检测与去除、信号标准化等环节，各环节相互配合形成完整的信号预处理流程：


（1）空间滤波技术用于处理多通道EEG的空间分布特性，通过利用不同电极间的空间关系来改善信号质量和提取特定的空间模式。共同平均参考是最常用的空间滤波方法，它通过从每个电极信号中减去所有电极的平均值，有效抑制全局共模噪声\cite{offner1950eeg}；

（2）时域滤波技术主要针对信号在时间维度上的噪声特性进行处理，是预处理流程中的基础环节。带通滤波器用于保留特定感兴趣的频段，如$\alpha$波（8--13Hz）、$\beta$波（13--30Hz）或$\gamma$波（30--50Hz），同时抑制其他频段的干扰信号。陷波滤波器则专门用于去除电源线干扰（50/60Hz），而高通滤波器主要用于消除低频基线漂移，提高信号的稳定性\cite{picton2000guidelines}；

（3）伪迹检测与去除技术主要针对采集过程中混入的各类生理和非生理伪迹进行识别和处理\cite{sweeney2012artifact}。传统的振幅阈值法通过设定绝对或相对阈值来检测突发性高幅伪迹，而小波分析方法能够在时频域同时定位异常活动。基于模板匹配的方法则适用于眼电和心电等具有典型波形特征的生理伪迹识别。现代伪迹处理方法更加强调保留而非简单剔除受污染数据，通过信号重建技术恢复潜在的神经活动模式。如图~\ref{fig:C1-5}所示，基于主成分分析的信号恢复、稀疏表示下的伪迹分离以及压缩感知理论下的信号重建等方法，在保留信号原有特性的同时，有效提升了信号质量和数据利用率；

（4）信号标准化技术通常作为预处理的收尾环节，旨在确保不同采集条件和被试间数据的可比性。基线校正通过减去或除以参考期的信号值，消除个体间基础活动水平的差异；幅度归一化则将不同量级的信号统一到相同尺度；电极位置标准化通过空间插值等技术，将不同电极布局的数据映射到标准化头皮坐标系，为跨实验数据整合和元分析提供支持。

\begin{figure}[htbp] % use float package if you want it here
  \centering
  \includegraphics[width=\textwidth]{C1-5ICA去伪迹.pdf}
  \caption{典型的基于ICA的去伪迹降噪算法效果\cite{albera2012ica}}
  \label{fig:C1-5}
\end{figure}
  
（二）特征提取和选择方法

特征提取是EEG分析的核心环节，旨在从预处理后的高维原始信号中提取具有判别力的低维特征表示。根据信号分析的侧重点不同，特征提取方法可分为时域特征、频域特征、时频特征、空间特征和非线性动力学特征等多个类别，各类方法从不同角度刻画神经信号的特性：

（1）时域特征提取方法直接从原始波形中获取特征信息，具有计算简单、物理意义明确的优势。统计特征如均值、方差、偏度和峰度描述了信号的基本分布特性，而Hjorth参数量化了信号的活动水平和复杂程度\cite{hjorth1970eeg}。过零率和峰值检测进一步反映了信号的振荡特性。时域特征常用于实时系统和快速筛选分析，适合ERP等具有明显时域特征的实验范式；
  
（2）频域特征提取方法通过功率谱分析来揭示信号的频率结构，是神经振荡研究的重要基础。快速傅里叶变换、自回归模型和周期图法是最常用的谱估计方法。基于这些方法可以提取各频段能量、频带功率比值、功率谱密度以及频谱熵等关键指标。频域特征在MI、精神负荷评估等任务中表现出色，能够有效捕捉$\alpha$、$\beta$节律的能量变化模式\cite{pfurtscheller1999event}；

（3）时频分析方法能够同时表征EEG的时间和频率动态特性\cite{cohen2019better}。短时傅里叶变换通过滑动窗口技术实现时频局部化，但其时频分辨率受限于海森堡不确定性原理。如图~\ref{fig:C1-6}所示，小波变换采用多分辨率分析策略，在不同时频尺度提供最优表示，适用于EEG中的瞬态成分。希尔伯特--黄变换则通过经验模态分解和希尔伯特谱分析，为高度非平稳信号提供自适应的时频表示。这些方法生成的时频图能够直观展示能量如何随时间和频率变化，对捕捉ERD/ERS等时变模式具有重要价值；

（4）空间特征提取方法关注多通道EEG的空间分布和电极间关系，能够反映大脑不同区域的功能联系。通道间相关性和相干性分析量化了不同脑区的信号相似度和同步性，而功能连接分析则构建了脑网络拓扑结构，揭示了信息流动和处理路径。共空间模式（Common Spatial Pattern, CSP）是目前最常用的空间特征提取方法之一 \cite{blankertz2008berlin}，它通过求解广义特征值问题，构造最能区分不同任务状态的空间滤波器，已被广泛应用于运动想象BCI系统，显著提高了分类性能；

（5）非线性动力学特征提取方法源于复杂系统理论，适用于捕捉神经系统的复杂动态特性。相关维数和Lyapunov指数描述了相空间中信号轨迹的几何特性和动力学稳定性，样本熵、排列熵等各类熵度量则量化了信号的复杂度和不规则性\cite{richman2000physiological}。近重复模式和递归定量分析进一步揭示了神经活动的周期性和确定性成分。这些特征能够表征神经信号的不同动态状态，在情绪识别、认知状态评估和神经疾病诊断中展现出较高的应用价值。

\begin{figure}[htbp] % use float package if you want it here
  \centering
  \includegraphics[width=\textwidth]{C1-6时频分析.pdf}
  \caption{基于连续小波变换的单通道EEG信号时频域分析}
  \label{fig:C1-6}
\end{figure}


在特征提取的基础上，特征选择构成了特征工程的重要补充环节，旨在从大量候选特征中筛选最具判别力的子集，同时降低维度并提高泛化能力\cite{guyon2003introduction}。信息增益、卡方检验等滤波法基于特征与目标变量的统计关系进行评分，计算简单但容易忽略特征间的相互作用。递归特征消除等包装法通过分类器性能来评估特征子集，效果优异但计算成本较高。基于L1正则化的嵌入式方法将选择过程融入模型训练，在效果和效率之间取得了良好平衡。合理的特征选择方法不仅能提高分类性能，还能降低计算复杂度，增强模型的可解释性。

（三）分类识别方法

分类识别算法是将提取的特征映射到目标类别的关键步骤，也是传统EEG分析流程的最终环节。根据算法原理和适用场景的不同，各种分类算法在EEG处理中发挥着不同的作用，需要根据具体应用需求进行选择和优化：

（1）线性判别分析（Linear Discriminant Analysis, LDA）因其计算效率高和良好的泛化能力，成为EEG分析特别是BCI系统中最常用的分类器之一\cite{lotte2010regularizing}。LDA基于各类别数据服从高斯分布且共享协方差矩阵的假设，通过构建线性决策边界来最大化类间分离。在特征维度不高且样本量有限的场景下，LDA的简单性和稳健性使其适合实时BCI系统。LDA的改进变体如正则化LDA和收缩LDA通过引入正则化项，进一步提高了在小样本高维数据上的性能表现；

（2）支持向量机（Support Vector Machine, SVM）通过寻找最大间隔超平面来分离不同类别，在EEG分类中也得到了广泛应用。核技巧使SVM能够有效处理线性不可分的复杂问题，其中径向基函数核在EEG分类中应用最为广泛\cite{rakotomamonjy2003variable}；

（3）K近邻（K-Nearest Neighbors, KNN）算法基于特征空间中的距离度量进行分类，实现简单直观且易于理解\cite{cover1967nearest}。该方法不需要显式训练过程，但在预测阶段需要计算查询样本与所有训练样本的距离，因此在大规模数据集上的效率相对较低；

（4）贝叶斯分类器基于概率模型进行决策，能够自然地量化分类的不确定性\cite{friedman1997bayesian}。朴素贝叶斯通过假设特征间条件独立，计算高效但表达能力相对有限。贝叶斯方法的突出优势在于能够整合先验知识，并在样本有限的情况下仍能给出合理的预测结果；

（5）决策树和随机森林方法通过构建逻辑决策规则进行分类，其结果具有良好的可解释性\cite{fraiwan2012automated}。单一决策树容易出现过拟合问题，而随机森林通过集成多棵在随机子空间训练的决策树，显著提高了泛化能力。


在实际应用中，EEG分类面临着多类分类和异步检测的挑战。现有方法多采用一对一、一对多等策略将二分类器扩展到多类问题，同时使用拒绝阈值和静息状态检测来实现异步控制。基于特征工程的EEG分析方法存在多方面的局限性。首先，特征设计主要依赖领域专家知识和人工经验，难以充分挖掘高维神经电数据中的隐含模式和复杂关系；其次，传统分析流程中的特征提取和分类环节通常采用分离处理方式。这种分段式处理方式可能导致信息损失和次优解，最优化的特征不一定是最佳分类特征，而最佳分类特征也可能不是最具生物学解释力的特征，这种矛盾在缺乏统一优化框架的情况下难以得到有效调和；第三，基于特征工程的方法本质上受限于传统的识别范式，难以支持开放式、创造性的脑机交互应用。现有方法通常将神经活动映射到有限的预定义类别或连续控制变量，这种固化的框架难以表达复杂的认知内容和意图的多样性，极大限制了BCI系统的表达能力。这些局限促使研究者积极探索更先进的EEG分析方法，特别是基于深度学习的表征学习方法和生成式建模方法。这些新兴方法有望克服特征工程的内在局限，实现更加灵活、鲁棒和表达力强的EEG处理系统。

\subsection{基于深度学习的脑电分析方法}
自2010年以来，基于深度神经网络的EEG分析方法在信号解码精度和应用范围方面都取得了显著突破。与传统基于特征工程的方法相比，深度学习方法展现出自动特征学习、端到端优化以及对复杂非线性关系的建模能力等核心优势，为解决EEG分析中的关键挑战提供了新思路。


如图~\ref{fig:C1-7}所示，卷积神经网络（Convolutional Neural Network, CNN）是应用最广泛的深度学习架构之一，其局部连接和权重共享特性特别适合处理具有空间结构的EEG信号\cite{liu2020epileptic}。DeepConvNet和ShallowConvNet针对EEG信号特性进行设计，直接从原始时间序列中学习空间滤波器和频率特征，在运动想象分类任务中超越了CSP结合LDA的方法\cite{schirrmeister2017deep}。为适应多通道EEG的空间拓扑结构，EEGNet设计专用的空间卷积层，首先学习通道间空间滤波器，随后学习时间卷积特征\cite{lawhern2018eegnet}。为捕捉多尺度时间特征，多分支CNN架构通过空洞卷积和多尺度卷积捕捉不同感受野的时间模式\cite{ang2012filter}。为联合建模特征，空-时-频多视图融合网络通过并行的时间、频率和空间卷积分支提取互补特征。
循环神经网络（Recurrent Neural Network, RNN）专门设计用于序列数据处理，能有效建模EEG的时间依赖性。在EEG分析中，RNN通常与CNN结合使用\cite{zhang2021survey}。CNN首先提取空间特征，RNN随后建模时间依赖性。这种CNN-RNN混合架构在情绪识别、癫痫检测等需要长期上下文的任务中表现优异。为增强时间建模能力，双向RNN、层级RNN等变体被提出，并通过课程学习等技术解决训练困难。

\begin{figure}[htbp] % use float package if you want it here
  \centering
  \includegraphics[width=\textwidth]{C1-7CNN.pdf}
  \caption{典型深度网络架构在脑电信号分析中的应用\cite{liu2020epileptic}}
  % \caption*{\centering (a) 卷积神经网络；(b) 循环神经网络。}
  \label{fig:C1-7}
\end{figure}


自编码器及其变体在EEG降噪、表征学习和异常检测中发挥重要作用\cite{vincent2008extracting}。这类模型通过“编码--解码”结构学习数据的低维表示，无需标签即可提取有用特征。自编码器被用于信号降噪和压缩表示，基于自编码器的异常检测方法也被应用于癫痫发作预测和脑部异常活动检测。
近年来，Transformer成为EEG分析的重要工具。其注意力机制通过学习动态权重，选择性关注信号中的关键成分，例如，空间注意力用于关注关键通道，时间注意力识别关键时间片段\cite{lee2022eeg}。基于Transformer的架构通过自注意力机制捕捉长距离依赖，在多种EEG分类任务中表现出色\cite{song2022eeg}。Vision Transformer的思想也被引入到EEG分析中，将时间序列划分为固定长度的时间片段作为“Token”，通过位置编码保持时序信息。多头注意力机制能够同时关注不同类型的模式，提供了比传统方法更丰富的特征表示能力。
图神经网络（Graph Neural Network，GNN）在处理EEG的空间拓扑结构方面展现出独特优势。通过将电极建模为图节点，电极间的空间或功能连接建模为图边，GNN能够有效捕获大脑网络的拓扑特性。图卷积网络（Graph Convolutional Networks，GCN）和图注意力网络（Graph Attention Networks，GAT）等方法在多种EEG任务中取得了良好效果，特别是在需要建模复杂空间关系的应用中表现突出。动态图神经网络还能够捕获时变的网络连接模式，为理解大脑动态网络提供了新的工具。

此外，深度学习可以与传统特征工程的方法相结合，通过结合时频特征的生理可解释性与深度网络的非线性表征能力，实现对复杂神经活动模式更精准的建模。多任务学习、迁移学习和元学习等技术也被广泛应用于EEG分析中，通过共享表示或快速适应机制提高模型的泛化能力和数据效率。

尽管基于深度学习的方法在EEG分析中取得了显著进展，但仍面临几个关键挑战：首先，现有模型设计往往缺乏神经生理机制的直接指导，多为通用深度学习架构的简单移植，难以充分利用EEG的特定结构和生物学特性；其次，这些方法通常需要大量标注数据才能有效训练，在EEG领域获取高质量标注数据成本高昂且耗时；第三，多模态融合能力有限，难以自然地整合EEG与其他生理或行为数据；最后，传统深度学习方法主要采用判别式范式，专注于将输入映射到预定义的类别或值，缺乏对EEG底层生成机制的建模能力。

EEG信号本质上是大脑神经网络复杂动力学过程的外在表现，其产生机制天然具有生成性特征。神经元群体的同步振荡、跨频段耦合以及时空动态演化等过程，都可以理解为一个复杂的概率生成过程。最近，深度学习领域本身正在经历从判别式模型向生成式模型的演进与融合。生成式建模作为深度学习的重要分支，恰好与EEG信号的这种内在生成性质相匹配，通过建模数据的概率分布，不仅能学习如何判断也能学习如何创造，这种能力有望实现更具表达力和生物合理性的EEG信号理解与应用系统。

\section{生成式建模方法研究现状}
\subsection{生成式建模的基本原理与发展}
生成式建模旨在学习观测数据的概率结构，从而实现采样、重建、插补与条件生成\cite{胡铭菲2022深度生成模型综述}。相较于仅刻画判别边界的判别式方法，生成式方法通过建模联合分布、边缘分布或条件分布，能够在密度估计、缺失值处理、数据增强、域迁移与多模态对齐等任务中提供统一的概率框架。形式上，生成式建模可分为两大范式：显式似然建模，典型代表包括自回归模型、变分方法与可逆流模型，通常可计算或可下界估计数据似然；隐式分布建模，典型代表包括对抗式方法与基于分数/扩散的方法，侧重直接生成高保真样本而非显式建模归一化的密度函数。

显式似然建模强调对数据似然或其下界的可计算性与可优化性，涵盖从经典概率模型到深度概率模型的连续谱系。经典模型如高斯混合模型（Gaussian Mixture Model, GMM）通过高斯分布加权叠加刻画复杂密度\cite{reynolds2009gaussian},隐马尔可夫模型（Hidden Markov Model, HMM）联合建模隐状态与观测序列，配合前向—后向与EM算法实现高效推断与学习\cite{rabiner2002tutorial, dempster1977maximum}。自回归模型通过链式法则将高维分布分解为条件分布乘积\cite{van2016conditional}，在语言、语音与图像像素层面密度建模中表现突出，但顺序采样通常较慢。
如图~\ref{fig:C1-8}所示，变分方法以变分自编码器（Variational Autoencoder, VAE）为代表\cite{kingma2013auto}，利用证据下界联合训练编码器与解码器学习连续、结构化的潜在表示。然而，常见解码器假设与优化目标可能导致样本细节偏平滑\cite{rezende2014stochastic}。
可逆流模型通过由可逆变换组成的映射将简单先验推送到数据分布，能够精确计算雅可比行列式并获得可评估的对数似然，典型结构包括耦合层与可逆卷积\cite{dinh2016density}、自回归流\cite{papamakarios2017masked}与连续正则化流\cite{chen2018neural}。该类方法提供可解释的密度与双向映射，但需在可逆性约束、参数效率与表达能力之间权衡，并关注自回归的采样速度或流模型的计算复杂度。

\begin{figure}[htbp] 
  \centering
  \includegraphics[width=\textwidth]{C1-8VAE.pdf}
  \caption{变分自编码器结构示意图\cite{胡铭菲2022深度生成模型综述}}
  \label{fig:C1-8}
\end{figure}

隐式分布建模侧重直接学习“如何生成”而非“如何归一化”，通过对比、分数或对抗目标驱动模型逼近数据流形。GAN以生成器--判别器博弈为核心\cite{goodfellow2014generative}，促使生成分布逐步贴近真实分布，衍生出深度卷积GAN\cite{radford2015unsupervised}、条件GAN\cite{hong2018conditional}、Wasserstein GAN\cite{arjovsky2017wasserstein}、以及面向高保真图像合成的StyleGAN\cite{karras2019style}、面向跨域映射的Pix2Pix\cite{isola2017image}与CycleGAN\cite{zhu2017unpaired}等。GAN在样本保真度方面具有优势，但易受训练不稳定、模式崩溃与评估困难等问题制约。基于分数与扩散的方法通过构造从数据到噪声的正向扰动过程与从噪声到数据的逆向去噪过程进行生成学习\cite{ho2020denoising,song2019generative}，可在离散或连续视角下训练，常采用去噪分数匹配等等价目标。潜在扩散进一步在压缩潜在空间内进行扩散\cite{rombach2022high}，结合跨模态条件（如文本条件的交叉注意力）显著提升效率与可控性。总体而言，隐式方法在高质量样本合成与多样性方面表现优异，但需通过正则化、架构设计与评估协议缓解训练不稳定与模式覆盖不足。

生成式模型的评估与范式和任务密切相关：显式似然模型可直接以对数似然或每维度比特数（Bits-per-dimension，BPD）等指标量化密度建模能力\cite{theis2015note}；隐式方法多依赖样本质量与多样性度量，例如初始分数（Inception Score，IS）\cite{salimans2016improved}和弗雷歇初始距离（Fréchet Inception Distance，FID）\cite{heusel2017gans}，以及面向下游任务的效用评估与人类主观评测\cite{zhang2018unreasonable}。随着应用扩展，安全性、公平性与知识产权等治理议题愈发重要，需在评估协议中纳入有害内容规避、偏见审计与内容溯源等维度。

\begin{figure}[htbp] 
  \centering
  \includegraphics[width=\textwidth]{C1-9LLM2.png}
  \caption{近年来大规模语言模型发展时间线\cite{zhao2023survey}}
  \label{fig:C1-9}
\end{figure}

在大规模预训练与基础模型的背景下，上述两大范式分别在不同模态与任务中取得系统化应用。自回归语言模型作为显式似然建模的代表，直接最大化条件似然驱动序列生成。如图~\ref{fig:C1-9}所示，大规模语言模型在近年来呈现出快速发展的趋势，从早期的GPT系列到后续的各类开源和闭源模型，模型规模和性能不断提升\cite{车万翔2023大模型时代的自然语言处理}；潜在扩散等隐式方法推动了文本—图像/视频合成与编辑的突破。此类技术通常结合大规模自监督/弱监督预训练、任务定制微调、多模态对齐与人类反馈强化学习（Reinforcement Learning from Human Feedback，RLHF）等对齐技术\cite{ouyang2022training}，以提升可控性、鲁棒性与符合人类偏好的程度。与此同时，数据治理、算力成本与安全对齐成为工程化落地的关键约束。


\subsection{生成式建模在生理时序信号中的应用}
生理时序信号呈现出强非平稳性、多尺度耦合与多源异质等统计属性，这要求生成式模型在时间依赖、频谱结构与跨通道相关性上同时刻画概率一致性\cite{赵雅琴2022基于改进深度生成对抗网络的心电信号重构算法}。显式似然家族（序列VAE、潜在连续动力学模型）通过可计算似然或其下界实现不确定性量化与校准，能够适配不规则采样、缺失与插补等场景；隐式分布家族（GAN及其稳健变体）与基于分数/扩散的方法则以对抗或分数目标直接逼近数据流形，具有较强的样本保真度与条件控制能力\cite{chung2015recurrent,rubanova2019latent,esteban2017real,yoon2019time,gulrajani2017improved,song2021score,tashiro2021csdi}。在这些范式之上，结构先验与约束的引入成为生理信号生成的关键设计要点。

围绕心电图（Electrocardiogram，ECG）等单/多导电生理序列，动力学合成模型提供了形态学与节律学的可控基线，并为深度模型的正则化与评测建立了参照\cite{mcsharry2003dynamical}。经典的ECG动力学模型通过微分方程系统描述心脏电活动的产生机制，能够生成具有生理学意义的合成心电信号。这些模型为理解生理信号的内在规律提供了重要基础，也为评估深度生成模型的生物学合理性提供了标准。
在隐式建模方面，改进的对抗训练与梯度惩罚缓解了训练不稳定与模式崩溃，使条件生成能够在形态与节律条件下保持分布覆盖与可控性\cite{gulrajani2017improved}，逼真度与统计一致性可以在对抗目标与结构性正则之间取得平衡\cite{thambawita2021deepfake}。时间序列GAN（TimeGAN）等方法通过同时优化监督损失和对抗损失，在保持时间动态的同时提高生成质量。条件GAN通过引入外部条件信息，能够实现对生成信号特定属性的控制。
面向缺失重建与不规则采样，带核过程先验的VAE类方法通过时间相关协方差建模实现插补的概率一致性\cite{fortuin2020gp}，而基于分数/扩散的条件模型则通过噪声--数据的双向过程学习多模态后验，兼顾长时间窗的依赖与覆盖度\cite{tashiro2021csdi}。高斯过程与神经网络的结合为处理不规则采样的生理信号提供了有效工具，能够在保持不确定性估计的同时实现高质量的插补。

与ECG相比，EEG的高维多通道与层级时频结构对生成式建模提出了更严格的先验契合需求。近期工作从“黑箱式合成”转向“先验增强”的范式。一方面，序列VAE与潜在连续动力学模型通过在潜在空间中编码慢变状态与快变观测，实现跨会话、跨段的带不确定性重建与插补，并可在不规则采样下保持后验的一致性\cite{chung2015recurrent,rubanova2019latent}；另一方面，隐式与混合式时序生成（如条件RNN-GAN与TimeGAN）通过对抗或自监督一致性目标耦合时间依赖与语义条件，同时配合谱域或时频约束以维持多频带统计与跨通道相干\cite{esteban2017real,yoon2019time}。进一步地，得分/扩散方法在多通道联合空间学习对数密度梯度，使联合插补与伪迹修复可在保持相干与相位结构的同时给出可校准区间，从而提升生成的不确定性可解释性\cite{song2021score,tashiro2021csdi}。

在生成模型的评价方面，逐渐形成以形态学可解释性、分布拟合质量与下游可用性为支点的三元框架，既关注波形、间期和振幅等可解释指标，也利用数据集标准化带来的统计可比性\cite{moody2001impact}。综上，生成式建模在EEG领域展现出从表征到推断的多重潜力。EEG生成式建模有望在保持生理可解释性的前提下实现更强的分布对齐与泛化能力，从而为多源EEG的生成、重建与跨模态对齐提供坚实基底。

\section{当前研究面临的问题与挑战}
EEG生成式建模作为人工智能与神经科学深度融合的前沿领域，在模型架构、算法优化与应用探索方面取得了初步进展。受限于EEG本身的复杂性、生成式建模技术的适应性局限以及脑机交互应用的高度特殊性，该领域仍面临一系列未被充分解决的科学问题。

（一）EEG内在复杂性与统一建模框架缺失

EEG作为大脑神经活动的直接电生理反映，其信号特性表现出极高的复杂性和多层次结构特征，而现有生成建模方法缺乏对这些内在复杂性的深度理解。

首先，EEG信号的多尺度时空动态特性对建模提出了根本性挑战。从时间维度看，EEG信号同时包含毫秒级的动作电位同步、秒级的认知事件相关电位、以及分钟级的认知状态变化等多个时间尺度的信息。这些不同时间尺度的神经活动相互嵌套和交互，形成了复杂的多层次时间结构。从空间维度看，EEG信号不仅反映局部神经元群体的活动，还受到远距离皮层连接、皮层--皮层下环路、以及体积传导效应等多种空间传播机制的影响。这种多尺度时空耦合特性使得传统的单一尺度建模方法难以充分捕获信号的完整信息，需要构建能够同时处理多个时空尺度的统一生成建模框架。

其次，EEG信号的非平稳性和个体差异性显著增加了建模复杂度。EEG信号的统计特性在时间上高度非平稳，受到被试的警觉水平、疲劳状态、情绪变化等多种因素的动态调制。同一被试在不同时间段记录的EEG信号可能表现出显著差异，这种时间非平稳性要求模型具备自适应的能力。同时，个体差异性表现在多个层面：解剖结构差异导致相同认知功能在不同个体间可能激活不同的脑区；头皮--皮层距离差异影响信号的幅度和空间分布；认知策略差异使得相同任务在不同个体间产生不同的神经响应模式。这些个体差异往往是非线性和高维的，传统的线性标准化方法难以有效处理。

最后，噪声与伪迹的复杂混合特性对信号质量构成严重威胁。EEG信号的信噪比通常较低，容易受到多种生理性和非生理性干扰。生理性伪迹包括眼电、肌电、心电等，这些伪迹的幅度往往远大于真实的脑电信号，且在时域和频域上与脑电信号高度重叠。非生理性干扰包括工频噪声、电极接触不良、运动伪迹等。不同类型的伪迹具有不同的产生机制、统计特性和时空分布模式，使得单一的去噪策略往往效果有限。更为复杂的是，某些伪迹成分可能与真实的神经信号在某些维度上相关，简单的滤波或去除可能导致有用信息的丢失。

（二）硬件约束下的信号质量优化挑战

EEG设备在实际应用中面临的硬件约束不仅限制了信号采集质量，更对整个系统的性能和可用性产生连锁影响，而现有技术缺乏系统性的解决方案。

首先，空间分辨率不足是EEG技术的根本性限制，其影响远超表面的采样密度问题。标准的高密度EEG系统虽然可以提供128-256个电极，但相对于大脑皮层约1,000平方厘米的表面积和数百亿神经元的复杂组织，这种采样密度仍然极其稀疏。空间分辨率不足导致的问题是多方面的：(1)空间混叠问题，相邻电极记录的信号高度相关，难以区分不同脑区的独立活动；(2)源定位精度低，逆问题求解的病态性使得从头皮电位重构皮层电流分布的准确性严重受限；(3)小范围脑区活动的检测能力不足，精细的认知功能往往涉及局部脑区的特异性激活，低空间分辨率难以捕获这些关键信息;(4)空间分辨率不足与个体差异相互作用，进一步加剧了跨个体泛化的困难。不同个体的头部形状、颅骨厚度、脑回形态等解剖差异使得相同电极位置记录到的信号来源可能完全不同，这种空间对应关系的不确定性严重影响了模型的泛化能力。

同时，便携性和实用性需求对EEG系统提出了严苛的工程约束。在移动健康监测、日常认知增强、便携式BCI等应用场景中，系统必须在尺寸、重量、功耗、成本等多个维度进行权衡优化。电极数量的减少不仅直接影响空间分辨率，还会导致冗余信息的丢失，降低系统对噪声和伪迹的鲁棒性。传统的通道选择方法主要基于统计显著性、信号质量或专家经验，缺乏对特定应用任务的针对性优化。更为关键的是，最优通道组合往往依赖于具体的认知任务、个体特征和环境条件，需要自适应的优化策略。此外，便携式设备还面临电极接触质量不稳定、运动伪迹增加、环境噪声复杂等挑战，要求系统具备更强的鲁棒性和自适应能力。


（三）无监督表征学习与跨域泛化技术瓶颈

EEG数据的标注稀缺性和跨域差异性使得传统监督学习范式面临根本性挑战，而现有无监督学习方法在EEG领域的适用性和有效性仍存在显著局限。

标注数据获取的根本性困难源于EEG数据标注的复杂性和专业性要求。与图像或文本数据不同，EEG信号的标注需要深厚的神经科学专业知识和昂贵的实验设备支持。高质量的EEG实验需要严格控制的实验环境、精密的同步系统、以及专业的被试招募和训练。更为重要的是，许多认知状态和神经活动模式本身就具有主观性和模糊性，即使是专家标注也存在较大的主观差异和不确定性。例如，在情绪识别任务中，同一段EEG信号可能被不同专家标注为不同的情绪类别；在认知负荷评估中，负荷水平的量化本身就存在争议。这种标注的不确定性不仅影响模型训练的质量，还使得模型评估变得困难。此外，标注过程的时间成本和人力成本都极为高昂，使得大规模标注数据集的构建在实践中几乎不可行。

跨域泛化困难是EEG分析领域的长期技术瓶颈，其根源在于EEG数据的多层次异质性。设备层面的差异包括电极数量、空间布局、采样频率、滤波参数等硬件规格的不同，以及放大器特性、阻抗要求、参考电极选择等技术细节的差异。这些设备差异直接影响信号的幅度、频谱特性和空间分布，使得在不同设备上记录的数据具有系统性偏差。实验协议的差异涉及任务设计、刺激材料、实验流程、环境控制等多个方面，这些差异会导致被试的认知策略、注意分配、疲劳水平等心理状态的不同，进而影响神经响应模式。被试群体的差异包括年龄、性别、教育背景、文化背景等人口学特征，以及认知能力、神经可塑性、疾病状态等生理心理特征。这些差异不仅影响基线脑活动水平，还会调制任务相关的神经响应。更为复杂的是，这些不同层面的域差异往往相互耦合，形成复合性的域偏移问题，使得简单的域适应方法难以有效处理。

现有自监督学习方法在EEG领域的适用性面临多重挑战。首先是预训练任务设计的困难，与图像领域的旋转、裁剪、颜色变换等数据增强方法不同，EEG信号的数据增强需要考虑生理合理性和神经科学约束。简单的时域截取、幅度缩放或噪声添加可能破坏信号的内在结构和生理意义。如何设计既能提供有效学习信号又能保持生理合理性的自监督任务，是一个尚未充分解决的问题。其次是掩码策略的复杂性，EEG信号的时空结构比自然图像更为复杂，简单的随机掩码可能无法充分利用信号的结构特性。需要考虑脑区功能分工、时频特征、以及个体差异等因素来设计更为精细的掩码策略。再次是模型架构的适配性问题，现有的预训练模型多借鉴计算机视觉或自然语言处理领域的架构，缺乏对EEG信号特殊性质的考虑。如何设计既能充分利用EEG信号特性又具备强大表征学习能力的模型架构，仍是一个开放性问题。

（四）系统集成与实用化部署的挑战

从理论算法到实用系统的转化过程中，生成式EEG建模技术面临着系统集成复杂度高、实时性要求严格、用户体验优化困难等多维度工程挑战。

认知状态建模的多维度复杂性是构建实用系统的核心难题。人类认知状态是一个多维度、动态变化的复杂系统，涉及注意力、记忆、情绪、疲劳、动机等多个心理维度的相互作用。传统的单一指标评估方法（如功率谱分析、事件相关电位等）往往只能反映认知状态的某个侧面，难以提供全面准确的状态刻画。更为复杂的是，不同认知维度之间存在复杂的相互作用和权衡关系。例如，高认知负荷可能伴随高专注度，但也可能导致注意力分散；疲劳状态可能降低认知效率，但在某些情况下也可能提高创造性思维。如何构建能够综合考虑多个认知维度及其相互关系的状态评价体系，实现对复杂认知状态的精准量化和动态追踪，是系统设计的根本性挑战。此外，认知状态的个体差异性和上下文依赖性进一步增加了建模复杂度，需要系统具备个性化适配和情境感知的能力。

多模态信息需求与技术集成的复杂性对系统架构提出了严苛要求。为了提高认知状态识别的准确性和鲁棒性，实用的认知监测系统需要提供多种生理指标的综合信息，而不仅限于EEG信号。例如，心率变异性可以反映自主神经系统的活动状态。这些不同模态的生理信号在反映认知状态时具有互补性，能够提供EEG信号无法单独捕获的认知信息。然而，集成多种传感器的系统在硬件设计、数据同步、信号质量控制等方面面临显著挑战。不同传感器的功耗、体积、佩戴舒适性要求各异，需要在系统集成度与用户体验之间寻求平衡。同时，多模态数据的实时处理和融合算法也需要针对计算资源受限的可穿戴设备进行优化设计。

用户交互设计与体验优化面临着技术可行性与用户需求之间的复杂平衡。有效的认知状态监测系统不仅需要准确的状态识别能力，还需要直观的状态呈现、合理的干预策略和良好的用户体验。状态可视化需要将复杂的多维认知状态以简洁直观的方式呈现给用户，既要保证信息的完整性又要避免认知负担。个性化建议生成需要结合用户的历史数据、当前状态、任务需求和个人偏好，提供具有针对性和可操作性的优化建议。自适应调节机制需要在不干扰用户正常活动的前提下，通过音频、视觉或其他感官通道提供适当的刺激和引导。这些功能的实现需要深入理解用户的认知特点、使用习惯和需求偏好，并通过迭代设计和用户测试不断优化。同时，还需要考虑隐私保护、数据安全、伦理规范等重要因素，确保系统的可信性和可接受性。


\section{主要研究内容与组织结构}
\subsection{本文的研究内容}
本文围绕多源EEG信号的生成式建模及应用开展研究，针对当前神经信号处理中普遍存在的信号结构复杂、跨域泛化能力不足以及系统化部署受限等关键问题，提出由信号级、表征级递进衔接的双层级生成式建模路径，形成从理论建模、方法实现到系统集成的解决方案。通过深入挖掘EEG信号的内在生成机制，建立从底层信号重构到高层语义理解的建模理论，为认知状态监测等脑机接口应用领域提供技术支撑。本文的主要研究内容如下：

（1）\textbf{多源EEG双层级生成式建模理论框架。}针对EEG信号的复杂时空结构和多尺度动态特性，本文构建信号级专注物理信号重建、表征级专注语义特征挖掘的分层生成建模架构，实现从底层信号重构到高层表征学习的统一建模范式。信号级层面将通过建模EEG信号的时域演化规律、空域传播机制和频域能量分布，解决数据质量受限、通道缺失与设备差异等基础性技术难题。该层级将学习EEG信号的内在生成机制，从低质量或不完整的观测中重构出高保真的神经活动模式，为后续分析提供可靠的数据基础。表征级层面将专注在无监督或弱监督条件下挖掘具有强泛化能力的潜在语义表征，从海量无标注数据中提取跨个体、跨任务的通用特征，有效缓解个体差异与域偏移带来的泛化困难。该框架通过形式化定义各层级的生成目标函数、结构化约束条件和层间信息传递机制，为多源EEG信号的理解提供统一的理论基础。

（2）\textbf{基于状态空间建模的EEG超分辨率方法。}针对EEG设备存在的电极数量有限、空间分辨率不足的硬件约束问题，本文将设计超分辨率方法，作为信号级生成式建模中的数据增强策略，实现从稀疏电极配置推断高密度EEG分布。拟设计基于状态空间模型的特征提取模块，有效捕获EEG信号的时序动态和潜在状态。将构建低分辨率特征提取器和高分辨率信号预测器的端到端级联生成架构，通过多尺度特征融合和渐进式上采样策略，实现从粗粒度到细粒度的信号重建过程。针对EEG信号的生理特性，将设计融合重建保真度损失、时域平滑性约束和空域连续性正则化的复合目标函数，确保生成信号在保持高保真度的同时具备良好的生理合理性。

（3）\textbf{协同引导学习驱动的EEG通道选择优化方法。}针对EEG设备在尺寸、重量、功耗和成本等方面的工程约束，本文将提出自适应通道选择框架，作为信号级生成式建模中的数据约简策略，实现在显著减少电极数量的前提下最大程度保持信号判别性能的智能优化。拟受认知神经科学中全局工作空间理论和模块化专门处理机制的启发，设计融合全局信息整合和任务特化处理的双路径协同引导学习架构。将开发联合卷积操作和状态空间动力学的核心模块，建模局部时空模式与长程依赖关系。拟构建端到端的可微分通道重要性评估和选择机制，自动学习和发现对特定认知任务最具判别价值的电极位置组合。

（4）\textbf{基于去噪掩码自编码的EEG表征学习方法。}针对EEG数据标注成本高昂和跨域泛化困难的问题，本文将构建大规模自监督预训练框架，作为表征级生成式建模的核心实现，实现从海量无标注数据中学习通用表征。拟提出基于脑区拓扑异质的数据划分方法，将均匀电极网格投影为反映大脑功能分区的层次化表征，有效解决不同设备间电极配置差异问题。将设计去噪伪标签生成器，通过重建型去噪任务为自监督学习提供稳健训练信号。拟采用非对称掩码自编码器结构，通过轻量化解码器迫使编码器学习紧凑的语义表征。该框架将在多个涵盖不同任务类型、设备配置和被试群体的公开数据集上进行预训练，通过零样本迁移和少样本微调验证跨域迁移的有效性。

（5）\textbf{认知状态监测与调节闭环系统原型。}基于上述双层级生成建模技术栈，本文将设计并实现认知状态监测与调节系统原型，验证生成式建模从理论到应用的完整转化路径。该系统将构建专注度--认知负荷二维评价体系，将认知状态空间划分为认知奋进、心流、游离、疲劳四个象限，实现认知状态的精准量化和可视化。拟集成多模态信息融合、大语言模型驱动的个性化建议生成、基于自适应音频的神经调制等功能模块，形成感知、分析和干预的闭环认知管理系统。通过系统原型的开发和测试，验证生成式人工智能技术在认知增强领域的应用潜力。

\subsection{本文的组织结构}
如图~\ref{fig:C1-10}所示，本文共分为七章，围绕多源EEG信号处理中信号质量受限、跨域泛化能力不足以及系统化部署困难等核心问题，构建了双层级生成式建模的完整研究体系。各章节内容相互支撑，从理论框架的建立，到信号级与表征级方法的深入探索，再到完整系统的工程实现，形成了从基础研究到应用验证的技术闭环。


\begin{figure}[htbp] % use float package if you want it here
  \centering
  \includegraphics[width=0.84\textwidth]{C1-10章节.pdf}
  \caption{本文主要研究内容和组织结构}
  \label{fig:C1-10}
\end{figure}

本章为绪论，阐述课题背景与意义，梳理相关工作与关键挑战，并概括了论文的研究内容。第二章提出双层级生成式建模理论框架，给出信号级、表征级的形式化定义、设计准则与层间耦合关系，为后续技术方法的开发奠定理论基础。第三章围绕信号级的超分辨率方法（Mamba-based Super-Resolution for EEG, MASER）进行详细论述，包括eMamba架构设计、训练目标与正则化策略以及重建质量与下游任务的综合评测，展示在信号增强方面的技术突破。第四章研究信号级的通道自适应选择方法（Cognitive State Space Model, CogSSM），介绍双路径共引导学习与特征提取模块设计，并开展跨数据集性能对比与通道可解释性分析，为便携式设备的优化部署提供技术支撑。第五章提出表征级自监督预训练（Denoising Masked AutoEncoder for EEG, DMAE-EEG）框架，详述脑区拓扑异质划分方法、去噪伪标签生成器与主干网络架构，在多任务与跨域迁移场景中进行验证，展现大规模预训练在EEG信号分析中的巨大潜力。第六章基于前述双层级生成式建模技术栈，完成智能认知增强头环原型的集成设计与功能验证，详述硬件平台实现、二维认知状态评价矩阵、系统原型评估与调节功能模块等，初步验证生成式建模技术在认知状态管理中的应用可行性，展现了从理论建模到原型应用的完整技术转化路径。第七章对全文进行总结，分析EEG信号生成式建模方法的局限，探讨多模态融合与大规模预训练的发展方向，并展望生成式人工智能与BCI深度融合的技术前景。

